\section{Capacitor}
An ideal capacitor is a two-terminal (one-port), passive, linear, time-invariant,
non-dissipative device that stores energy (\& thus charge) in the
form of an electric field.

\begin{itemize}
  \item \textbf{Capacitance:}
    The relationship between charge and
    voltage is expressed by a proportionality factor called capacitance
    (\(C\)) and is measured in Farads (F). It also gives us the
    capacitor's version of Ohm's law:
    \begin{align*}
      Q &= C \cdot V \\
      \text{Differentiating both sides with respect to time:} \\
      \frac{dQ}{dt} &= C \cdot \frac{dV}{dt} \hspace{1cm} \text{(assuming \(C\) is constant)} \\
      I &= C \cdot \frac{dV}{dt} \hspace{1cm} [I = \frac{dQ}{dt}]
    \end{align*}
    The inverse of capacitance is called elastance (\(S\)) and is measured in Farads
    inverse (F\(^{-1}\)).
  \item \textbf{Impedance:}
    Unlike resistor where a proportionality factor directly relates
    voltage and current, a capacitor is different, the relationship
    between voltage and current is related instead by a differential
    equation. We thus can not describe resistance of a capacitor in the
    same way we do for a resistor. We must somehow trick the above
    capacitor relationship into giving us a proportionality factor
    between voltage and current. Laplace Transform can be used to
    achieve this,
    \begin{align*}
      I &= C \cdot \frac{dV}{dt} \\
      I(s) &= C \left( s V(s) - V(0) \right) \hspace{1cm} [\mathcal{L}\{\tfrac{dV}{dt}\} = sV(s) - V(0)] \\
      I(s) &= sC\, V(s) \hspace{1cm} [\text{assuming \(V(0) = 0\)}] \\
      V(s) &= \frac{1}{sC} I(s)
    \end{align*}
    We can now define the impedance of a capacitor as the ratio
    of voltage to current in the Laplace domain:
    \[
      Z_C = \frac{V(s)}{I(s)} = \frac{1}{s C} = \frac{1}{j\omega C} = \frac{1}{j 2 \pi f C} \hspace{1cm} [s = j\omega = j 2\pi f]
    \]
  \item \textbf{Admittance:}
    The inverse of impedance is called admittance (\(Y\)) and
    is measured in Siemens (S). The relationship between impedance
    and admittance is:
    \[
      Y = \frac{1}{Z} = \frac{1}{\frac{1}{sC}} = sC
    \]

    You'll notice when we apply KCL to circuits within this book,
    we'll mostly use admittance \(Y\) to express the impedance, we
    do this as it simplifies the analysis a lot. The admittance
    version of Ohm's law is:
    \[
      I = V \cdot Y
    \]
  \item \textbf{Frequency Response:}
    The impedance of a capacitor is inversely proportional to
    frequency. This means that at low frequencies, the impedance
    of a capacitor is very high and at high frequencies, the impedance
    of a capacitor is very low. This is why at high frequencies,
    the capacitor starts to load the circuit which is often a
    very desirable effect in many applications. This is also why
    capacitors are used for decoupling (they act as a short
    circuit at high frequencies).
  \item \textbf{Energy and Power:}
    The energy and power stored in a capacitor is given by:
    \begin{align*}
      I &= C \cdot \frac{dV}{dt} \\
      P &= V \cdot I \\
      P &= V \cdot C \cdot \frac{dV}{dt} \\
      P &= C \cdot V \cdot \frac{dV}{dt} \\
      P &= C \cdot \frac{1}{2} \cdot \frac{d(V^2)}{dt} \hspace{1cm} [\frac{d(V^2)}{dt} = 2V \cdot \frac{dV}{dt}] \\
      E &= \int_0^t P \, dt' = \int_0^t C \cdot \frac{1}{2} \cdot \frac{d(V^2)}{dt'} \, dt' \\
      E &= \frac{1}{2} C \left[ V^2(t) - V^2(0) \right] \\
      E &= \frac{1}{2} C V^2 \hspace{1cm} [V(0) = 0,\ V = V(t)]
    \end{align*}
  \item \textbf{Types of Capacitors}:
    There are many types of capacitors, some of the most common are:
        \begin{itemize}
      \item \textbf{Ceramic/MLCC Capacitor}: Small size, low cost \& a
        very low equivalent series resistance (ESR) and inductance (ESL) make
        them really good for high frequency applications. The downside
        is that \textbf{capacitance varies greatly with applied DC voltage} (a plus
        if you want to use them as cheap varactors) and it also
        \textbf{varies greatly with temperature}. These usually have \textbf{low
        capacitance values} for a given size. Ceramic capacitors are divided
        into classes,
        \begin{itemize}
          \item Class 1: C0G/NP0 (high stability, low capacitance)
          \item Class 2: X7R (moderate stability, moderate capacitance)
          \item Class 3: Y5V (low stability, high capacitance)
        \end{itemize}
        Examples: 100 nF X7R 0603 ceramic capacitor across a supply pin, or
        a C0G/NP0 capacitor in a precision timing or filter network.
      \item \textbf{Electrolytic Capacitor}: High capacitance for a
        given case size, low cost and good for bulk energy
        storage. These usually have a \textbf{higher ESR and ESL, which
        translates to poor high frequency behaviour}, and are \textbf{almost
        always polarized} (reverse voltage may destroy the part). These
        also have a \textbf{finite lifetime} (electrolyte evaporation, sped up
        by heat and ripple current).
        Example: 470\,\(\mu\)F aluminum electrolytic capacitor on the output
        of a linear or switching power supply.
      \item \textbf{Tantalum Capacitor}: Higher capacitance density
        than ceramics and more stable over voltage and temperature
        than electrolytics. These also have a lower ESR/ESL than
        electrolytics. They are very \textbf{often polarized} (known for
        failing short and igniting under fault) and \textbf{derating (keeping
        applied voltage well below rated) is strongly recommended} if
        you want reliability. These are often \textbf{expensive}. Example:
        10\,\(\mu\)F tantalum capacitor used for bulk decoupling on a densely
        populated board where space is at a premium.
      \item \textbf{Film Capacitor}: Excellent stability over temperature
        and voltage, very low dielectric absorption, and low ESR/ESL.
        These are non-polarized but have \textbf{a larger and more expensive per
        unit capacitance} than other types. They are also \textbf{limited to
        lower capacitance values}. Usually used in high quality audio,
        precision timing and filters. Example: 1\,\(\mu\)F polypropylene film
        capacitor in a high quality audio crossover network.
    \end{itemize}
  \item \textbf{They can form everywhere}:
    A capacitor exists anywhere where you have two conductors
    separated by an insulator. The charges are stored on the
    surface of the conductors and the energy is stored in the
    electric field between the conductors.
  \item \textbf{Circuit Analysis:}
    The behaviour of capacitor is affected by the kind of Input Signal (AC
    or DC) applied to it. Find the admittance \(sC\) in the Small Signal,
    Frequency Domain, AC Input Analysis (Q4) before deciding whether to
    Short, Open or keep the capacitor as is within your analysis. Table~\ref{tab:hand-analysis-simplifications}
    serves as a quick reference.
\end{itemize}
